\documentclass{rapport}
\usepackage[utf8]{inputenc}
\usepackage{lipsum}
\usepackage{dirtytalk}
\usepackage[style=ieee]{biblatex} %Imports biblatex package
\usepackage{csquotes}
\usepackage{amsmath}
\usepackage{tcolorbox}
\usepackage{pdfpages}
\usepackage{outlines}
\usepackage{multirow}
\usepackage{xcolor}
\usepackage{colortbl}
\usepackage[T1]{fontenc}
\usepackage{xcolor}
\usepackage{listingsutf8}
\usepackage{anyfontsize}
\usepackage{enumitem}
\usepackage{graphicx}
\usepackage{hyperref}
\usepackage{array}
\usepackage{tikz}
\usepackage[para]{footmisc}
\usepackage{caption}
\usepackage{enumitem}
\usepackage{multicol}
%\usepackage{tocloft}

\setlist[itemize,1]{label=$\bullet$}

\hypersetup{
    colorlinks=true,
    linkcolor=blue,
    filecolor=magenta,
    urlcolor=cyan,
    }
\urlstyle{same}

\newlength\tindent
\setlength{\tindent}{\parindent}
\setlength{\parindent}{0pt}
\renewcommand{\indent}{\hspace*{\tindent}}

% custom commands
\newcommand{\newparagraph}{\vspace{0.5em}}
\newcommand{\itembullet}{\item[$\bullet$]}

\addbibresource{sources.bib} % premier fichier bilbiographique

\title{LSA_Lab_01_Template_Rapport_Première_Manip} %Titre du fichier

\definecolor{backgroundcolor}{rgb}{0.2, 0.2, 0.2} % Gris foncé, plus lisible
\definecolor{textcolor}{rgb}{1.0, 1.0, 1.0} % Blanc
\definecolor{commentcolor}{rgb}{1.0, 0.6, 0.0} % Orange clair
\definecolor{keywordcolor}{rgb}{0.5, 0.5, 1.0} % Bleu foncé
\definecolor{stringcolor}{rgb}{0.8, 0.4, 0.4} % Rouge foncé


% Configurer listings
\lstset{
  language=bash,
  inputencoding=utf8/latin1,
  extendedchars=true,
  backgroundcolor=\color{backgroundcolor}, % Couleur de fond
  basicstyle=\ttfamily\small\color{textcolor}, % Couleur du texte
  commentstyle=\color{commentcolor}, % Couleur des commentaires
  keywordstyle=\color{keywordcolor}, % Couleur des mots-clés
  stringstyle=\color{stringcolor}, % Couleur des chaînes de caractères
  numbers=left,
  numberstyle=\tiny\color{gray},
  stepnumber=1,
  frame=single,
  breaklines=true,
  showstringspaces=false,
  captionpos=b % Caption en bas
  literate={à}{{\`a}}1 {é}{{\'e}}1 {è}{{\`e}}1 {ê}{{\^e}}1 {ë}{{\¨e}}1 {â}{{\^a}}1 {î}{{\^i}}1 {ï}{{\¨i}}1 {ô}{{\^o}}1 {ù}{{\`u}}1 {û}{{\^u}}1 {ü}{{\¨u}}1 {ç}{{\c{c}}}1
}

\setlist[enumerate,1]{label=\arabic*., ref=\arabic*}
\setlist[enumerate,2]{label=\alph*., ref=\alph*} % Change le label pour les sous-points

\begin{document}

%----------- Informations du rapport ---------

\logo{img/henallux.jpg}
\dep{MASI}
\cours{Analyse programmation}
\titre{Rapport d'analyse programmation} %Titre du fichier .pdf
\sujet{Rapport du cours d'Analyse programmation} %Nom du sujet

\enseignant{\textsc{Huygens} Adrien} %Nom de l'enseignant

\eleves{\textsc{Dujardin} Xavier \\
        \textsc{Devrient} Félix \\
        \textsc{Lerat} Kylian \\
        \textsc{Croibien} Robin \\
        \textsc{Kinet} Rémi \\
        \textsc{Monis} Lucas}


%----------- Initialisation -------------------

\fairemarges%Afficher les marges
\fairepagedegarde%Créer la page de garde

\pagenumbering{gobble}

\hbox{}
\newpage

\tabledematieres

\pagenumbering{arabic}
\setcounter{page}{1}

\section{Introduction}

\newpage


\phantomsection
\printbibliography[heading=bibintoc, title={Bibliographie}]

%\newpage

%\phantomsection
%\addcontentsline{toc}{section}{Table des figures}
%\listoffigures
%\newpage

%\phantomsection
%\addcontentsline{toc}{section}{Table des Codes}
%\section*{Table des Codes}
%\lstlistoflistings


\end{document}